% ------------------------------------------------------------
% Chapter 9: Large-Scale Structure
% ------------------------------------------------------------

\chapter{Large--Scale Structure}
\label{chap:lss}

\section{Overview}

Large--scale structure (LSS) measurements probe the distribution of matter
at cosmological scales. In \rsvpabbr{}, the growth of structure is governed
by both gravitational and lamphrodynamic effects, leading to potentially
distinct predictions for the matter power spectrum, baryon acoustic
oscillations (BAO), and weak lensing.

\section{Matter Power Spectrum}

The matter power spectrum $P(k)$ encodes the statistical distribution of
density fluctuations as a function of wavenumber $k$. In \rsvpabbr{}, one
obtains modified linear growth factors and transfer functions, resulting in
\[
  P_{\mathrm{RSVP}}(k)
   = |T_{\mathrm{RSVP}}(k)|^2 P_{\mathrm{prim}}(k).
\]
The transfer function $T_{\mathrm{RSVP}}(k)$ depends on the lamphrodynamic
couplings and can differ from the \LCDM{} transfer function, especially on
large scales.

\section{Baryon Acoustic Oscillations}

BAO provide a standard ruler for cosmological distances. Lamphrodynamic
effects alter both the sound horizon and the growth of baryonic
oscillations, potentially shifting the BAO peak. Comparing BAO in \rsvpabbr{}
with galaxy surveys constrains deviations from \LCDM{}.

\section{Weak Lensing}

Weak lensing measurements probe the projected mass distribution along the
line of sight. In \rsvpabbr{}, lamphrodynamic density contributes to the
convergence $\kappa$, as in \cref{chap:dark-sector}. Combining lensing with
BAO and galaxy clustering improves constraints on $\rho_{\mathrm{lam}}$ and
$p_{\mathrm{lam}}$.

\section{Discussion}

Large--scale structure provides important constraints on \rsvpabbr{}, though
significant degeneracies with standard cosmological parameters remain.
Future surveys may help break these degeneracies and provide stronger tests
of lamphrodynamic cosmology.

\clearpage
